%%%%%%%%%%  scan idx 142  |  folio 131  %%%%%%%%%%

\par\medskip
\noindent\textbf{Definición 6.6.}
Sea $X$ un $k$-esquema de tipo finito. Un punto cerrado $z\in X$ se
denomina una \emph{singularidad cuadrática ordinaria racional sobre
$k$} si se satisfacen las condiciones siguientes.

\smallskip
\noindent\textbf{(i)} Si $\operatorname{char}k\neq2$, o bien
$\operatorname{char}k=2$ y $\dim\mathcal O_{X,z}$ es impar, entonces
$z$ es una singularidad cuadrática no degenerada racional sobre $k$.

\smallskip
\noindent\textbf{(ii)} Si $\operatorname{char}k=2$ y
$\dim\mathcal O_{X,z}$ es par, entonces
\[
\overline k\otimes_k\widehat{\mathcal O}_{X,z}
 \simeq
\overline k[[X_0,X_1,\ldots,X_{2n}]]/(f),
\]
donde $\overline k$ es una clausura algebraica de $k$ y
\[
f=X_0^2+X_1X_2+X_3X_4+\cdots+X_{2n-1}X_{2n}.
\]

\par\medskip
\noindent\textbf{Lema 6.7.}
Sea $k$ un cuerpo de característica $2$, y pongamos
\[
A_0=k[[X_0,X_1,\ldots,X_{2n}]]/(f),
\qquad
f=X_0^2+X_1X_2+X_3X_4+\cdots+X_{2n-1}X_{2n}.
\]
Sea $(R,\mathfrak X)$ una deformación versal formal de
$X_0=\operatorname{Spec}(A_0)$. Entonces
\[
R=\Lambda[[t,t_1]],
\]
donde $t,t_1$ son generadores formales libres de $R$ sobre
$\Lambda$, y
\[
I_R(\mathfrak X|R)=(t^2).
\]
Además, $\mathfrak X\longrightarrow\operatorname{Spec}(\Lambda)$ es
formalmente liso.

\smallskip
\noindent\textit{Demostración.}
Escribamos $\mathfrak X=\operatorname{Spf}(A)$. Como
\[
\left(f,\frac{\partial f}{\partial X_0},\ldots,
         \frac{\partial f}{\partial X_{2n}}\right)
=(X_0^2,X_1,\ldots,X_{2n}),
\]
tenemos
\[
D^1(A_0|k,A_0)\simeq k[X_0]/(X_0^2).
\]
Así, $[\underline M_{X_0}(k[\varepsilon])]$ es un espacio vectorial
sobre $k$ de dimensión dos. Por 4.14,
\[
R\simeq\Lambda[[t,t_1]],
\qquad
A\simeq R[[X_0,X_1,\ldots,X_{2n}]]/(F),
\]
donde
\[
F=f+t_1+tX_0.
\]
La sucesión exacta
\[
0\longrightarrow A\xrightarrow{\ dF\ }
A\,dX_0\oplus\cdots\oplus A\,dX_{2n}
\longrightarrow\widehat\Omega_{A|R}\longrightarrow0,
\]
en la que
\[
dF:\ 1\longmapsto
\sum_{0\leq i\leq2n}\frac{\partial f}{\partial X_i}\,dX_i
\sisi{}{+t\,dX_0},
\]
da
\[
I_A(A|R)
=(t,X_1,\ldots,X_{2n})
=(X_0^2+t_1,t,X_1,\ldots,X_{2n}).
\]
% La fórmula impresa omite el sumando t\,dX_0, aunque
% F=f+t_1+tX_0 y el ideal siguiente fuerzan este término diferencial.

%%%%%%%%%%  scan idx 143  |  folio 132  %%%%%%%%%%

Por tanto,
\[
\begin{aligned}
A/I_A(A|R)
&\simeq A/(X_0^2+t_1,t,X_1,\ldots,X_{2n})\\
&\simeq\Lambda[[t_1,X_0]]/(X_0^2+t_1).
\end{aligned}
\]
Como $R$-módulo, este último módulo es isomorfo a
$R/(t)\oplus R/(t)$. Se sigue que
\[
I_R(A|R)=(t^2),
\]
lo cual demuestra la primera afirmación. Por otra parte,
\[
\begin{aligned}
A
&\simeq
\Lambda[[t,t_1,X_0,X_1,\ldots,X_{2n}]]/
 (X_0^2+X_1X_2+\cdots+X_{2n-1}X_{2n}+tX_0+t_1)\\
&\simeq\Lambda[[t,X_0,X_1,\ldots,X_{2n}]],
\end{aligned}
\]
de modo que $A$ es formalmente liso sobre $\Lambda$.
\hfill$\square$

El Lema 6.7 da el siguiente análogo del Teorema 6.4; dejamos su
demostración al lector.

\par\medskip
\noindent\textbf{Corolario 6.8.}
Sea $X_0$ un $k$-esquema de tipo finito cuyos puntos no lisos sean
aislados, supongamos que
\[
H^2(X_0,\underline D^0_{X_0})=0,
\]
y sea $(R,\mathfrak X)$ una deformación versal formal de $X_0$.
Denotemos por $\{z_1,\ldots,z_m\}$ el conjunto de puntos no lisos.
Supongamos que $\operatorname{char}k=2$, que $\dim X_0$ es par y que
todos los $z_i$ son singularidades cuadráticas ordinarias. Entonces:

\smallskip
\noindent\textbf{(1)} $R$ es formalmente liso sobre $\Lambda$.

\smallskip
\noindent\textbf{(2)} Cada ideal
$I_R(\mathfrak X_{(z_i)}|R)$ es principal, al igual que
\[
I_R(\mathfrak X|R)
=\prod_{1\leq i\leq m}I_R(\mathfrak X_{(z_i)}|R).
\]
Más precisamente,
\[
I_R(\mathfrak X_{(z_i)}|R)=(t_i^2),
\]
y $t_1,\ldots,t_m$ pueden completarse hasta un sistema de generadores
formales libres de $R$ sobre $\Lambda$.

\smallskip
\noindent\textbf{(3)} El morfismo
$\mathfrak X\longrightarrow\operatorname{Spec}(\Lambda)$ es
formalmente liso.

\smallskip
\noindent\textbf{(4)} Si $k$ es algebraicamente cerrado y $X_0$ es
una curva completa irreducible, entonces
\[
\dim\underline M_{X_0}=3g-3+a,
\]
donde
\[
g=\dim_kH^1(X_0,\mathcal O_{X_0}),
\qquad
a=\dim_kH^0(X_0,\underline D^0_{X_0}).
\]
